% Appendix: full tip index
\chapter{Tips Index\quad TIP 总索引}

\begin{longtable}{@{}r p{0.44\textwidth} p{0.40\textwidth}@{}}
\toprule
\textbf{\#} & \textbf{English} & \textbf{中文} \\
\midrule
\endhead
1 & Care About Your Craft & 关心你的手艺 \\
2 & Think! About Your Work & 思考！关于你的工作 \\
3 & Provide Options, Don't Make Lame Excuses & 提供选项，别找蹩脚的借口 \\
4 & Don't Live with Broken Windows & 不要容忍破窗 \\
5 & Be a Catalyst for Change & 做变革的催化剂 \\
6 & Remember the Big Picture & 记住大局 \\
7 & Make Quality a Requirements Issue & 把质量当作需求问题 \\
8 & Invest Regularly in Your Knowledge Portfolio & 定期投资你的知识资产 \\
9 & Critically Analyze What You Read and Hear & 批判性地分析你所读所闻 \\
10 & It's Both What You Say and the Way You Say It & 说什么与怎么说同样重要 \\
11 & DRY - Don't Repeat Yourself & DRY——不要重复自己 \\
12 & Make It Easy to Reuse & 让复用变得容易 \\
13 & Eliminate Effects Between Unrelated Things & 消除不相关事物之间的影响 \\
14 & There Are No Final Decisions & 没有最终决定 \\
15 & Use Tracer Bullets to Find the Target & 用曳光弹寻找目标 \\
16 & Prototype to Learn & 用原型学习 \\
17 & Program Close to the Problem Domain & 贴近问题领域编程 \\
18 & Estimate to Avoid Surprises & 估算以避免意外 \\
19 & Iterate the Schedule with the Code & 让进度表随代码迭代 \\
20 & Keep Knowledge in Plain Text & 让知识保存在纯文本中 \\
21 & Use the Power of Command Shells & 善用命令行的力量 \\
22 & Use a Single Editor Well & 熟练使用一个编辑器 \\
23 & Always Use Source Code Control & 始终使用源码版本控制 \\
24 & Fix the Problem, Not the Blame & 解决问题，而不是归咎于人 \\
25 & Don't Panic & 不要惊慌 \\
26 & ``select`` Isn't Broken & ``select'' 没有坏 \\
27 & Don't Assume It - Prove It & 别假设——去证明 \\
28 & Learn a Text Manipulation Language & 学一门文本处理语言 \\
29 & Write Code That Writes Code & 写能生成代码的代码 \\
30 & You Can't Write Perfect Software & 你写不出完美的软件 \\
31 & Design with Contracts & 用契约来设计 \\
32 & Crash Early & 尽早崩溃 \\
33 & If It Can't Happen, Use Assertions to Ensure That It Won't & 若不可能发生，就用断言确保它不会 \\
34 & Use Exceptions for Exceptional Problems & 用异常处理异常问题 \\
35 & Finish What You Start & 有始有终 \\
36 & Minimize Coupling Between Modules & 最小化模块间的耦合 \\
37 & Configure, Don't Integrate & 配置，而非集成 \\
38 & Put Abstractions in Code, Details in Metadata & 抽象进代码，细节进元数据 \\
39 & Analyze Workflow to Improve Concurrency & 分析工作流以改善并发 \\
40 & Design Using Services & 用服务来设计 \\
41 & Always Design for Concurrency & 始终为并发而设计 \\
42 & Separate Views from Models & 让视图与模型分离 \\
43 & Use Blackboards to Coordinate Workflow & 用黑板协调工作流 \\
44 & Don't Program by Coincidence & 不要靠巧合编程 \\
45 & Estimate the Order of Your Algorithms & 估算算法的量级 \\
46 & Test Your Estimates & 检验你的估算 \\
47 & Refactor Early, Refactor Often & 尽早重构，经常重构 \\
48 & Design to Test & 为测试而设计 \\
49 & Test Your Software, or Your Users Will & 你不测试，用户就会替你测试 \\
50 & Don't Use Wizard Code You Don't Understand & 不要用你不理解的向导代码 \\
51 & Don't Gather Requirements - Dig for Them & 不要收集需求——去挖掘需求 \\
52 & Work with a User to Think Like a User & 与用户共事，像用户一样思考 \\
53 & Abstractions Live Longer than Details & 抽象比细节活得更久 \\
54 & Use a Project Glossary & 使用项目术语表 \\
55 & Don't Think Outside the Box - Find the Box & 别跳出框框想——先找到框框 \\
56 & Listen to Nagging Doubts - Start When You're Ready & 倾听内心的疑虑——准备好了就开始 \\
57 & Some Things Are Better Done than Described & 有些事做胜过说 \\
58 & Don't Be a Slave to Formal Methods & 不要做形式化方法的奴隶 \\
59 & Expensive Tools Do Not Produce Better Designs & 昂贵的工具不产生更好的设计 \\
60 & Organize Around Functionality, Not Job Functions & 围绕功能组织，而非职务 \\
61 & Don't Use Manual Procedures & 不要用手工流程 \\
62 & Test Early. Test Often. Test Automatically. & 早测试，常测试，自动化测试 \\
63 & Coding Ain't Done 'Til All the Tests Run & 所有测试都跑通，编码才算完 \\
64 & Use Saboteurs to Test Your Testing & 用“破坏者”来测试你的测试 \\
65 & Test State Coverage, Not Code Coverage & 测试状态覆盖，而非代码覆盖 \\
66 & Find Bugs Once & 同一个 bug 只找一次 \\
67 & Treat English as Just Another Programming Language & 把英语当成另一种编程语言 \\
68 & Build Documentation In, Don't Bolt It On & 把文档做进去，而不是补上去 \\
69 & Gently Exceed Your Users' Expectations & 温和地超越用户的期望 \\
70 & Sign Your Work & 为你的作品署名 \\
\bottomrule
\end{longtable}
